\documentclass[preview]{standalone}

\usepackage{amsmath}
\usepackage{amssymb}
\usepackage{stellar}
\usepackage{definitions}

\begin{document}

\id{polynomial-rings-exercises}
\genpage

\section{Exercises}

\begin{snippetexercise}{polynomial-quotient-units-zero-divisors-exercise}{Units and zero divisors in a polynomial quotient}
    Let
    \[
        f=x^3-2x^2-2x-3\in\rationalnumbers[x],
        \qquad
        I=(f),
        \qquad
        A=\rationalnumbers[x]/I.
    \]
    Determine whether
    \[
        (x^2-x+1)+I
        \qquad\text{and}\qquad
        (x^2-2x-3)+I
    \]
    are units or zero divisors in \(A\).
\end{snippetexercise}

\begin{snippetsolution}{polynomial-quotient-units-zero-divisors-solution}{Units and zero divisors in a polynomial quotient}
    Set
    \[
        g_1=x^2-x+1,
        \qquad
        g_2=x^2-2x-3.
    \]
    Since
    \[
        f=(x-3)(x^2+x+1),
        \qquad
        g_2=(x-3)(x+1),
    \]
    one has
    \[
        g_2(x^2+x+1)=(x+1)f.
    \]
    Passing to the quotient gives
    \[
        (g_2+I)\bigl((x^2+x+1)+I\bigr)=0_A.
    \]
    Both factors are nonzero because their representatives are nonzero
    and have degree smaller than \(\polynomialdeg f=3\). Thus \(g_2+I\)
    is a zero divisor.

    The polynomial \(g_1\) is irreducible over \(\rationalnumbers\) and
    does not divide \(f\), so \(\gcd(f,g_1)=1\). The Euclidean algorithm
    gives
    \[
        f=(x-1)g_1+(-4x-2)
    \]
    and
    \[
        g_1
        =
        \left(-\frac{x}{4}+\frac{3}{8}\right)(-4x-2)
        +\frac{7}{4}.
    \]
    Put \(q_0=x-1\) and
    \(q_1=-x/4+3/8\). Back-substitution yields
    \[
        1
        =
        \frac{4}{7}(1+q_1q_0)g_1-\frac{4}{7}q_1f.
    \]
    Therefore \(g_1+I\) is a unit, with inverse
    \[
        (g_1+I)^{-1}
        =
        \left(
            -\frac{1}{7}x^2+\frac{5}{14}x+\frac{5}{14}
        \right)+I.
    \]
    In particular, \(A\) is not an \integraldomain, equivalently
    \((f)\) is not a prime \ideal of \(\rationalnumbers[x]\).
\end{snippetsolution}

\begin{snippetexercise}{polynomial-quotient-field-units-remainder-exercise}{A reducible polynomial quotient}
    Let
    \[
        f=X^3-2X^2-2X+4\in\rationalnumbers[X],
        \qquad
        A=\rationalnumbers[X]/(f),
    \]
    and put
    \[
        g_1=X^2+4X+4,
        \qquad
        g_2=4X-2X^2.
    \]
    Determine:
    \begin{enumerate}
        \item whether \(A\) is a field or an integral domain;
        \item whether \(g_1+(f)\) and \(g_2+(f)\) are units;
        \item a representative of degree less than \(3\) for
        \(X^4+(f)\).
    \end{enumerate}
\end{snippetexercise}

\begin{snippetsolution}{polynomial-quotient-field-units-remainder-solution}{A reducible polynomial quotient}
    One has
    \[
        f=(X-2)(X^2-2).
    \]
    Thus \((f)\) is not prime, and \(A\) is neither an
    \integraldomain nor a field. Explicitly, the two nonzero classes
    of \(X-2\) and \(X^2-2\) have zero product.

    Since
    \[
        g_2=-2X(X-2),
        \qquad
        g_2(X^2-2)=-2Xf,
    \]
    the nonzero class \(g_2+(f)\) is a zero divisor and is not
    invertible.

    On the other hand, \(g_1=(X+2)^2\) is coprime to \(f\). The
    Euclidean divisions
    \[
        f=(X-6)g_1+(18X+28)
    \]
    and
    \[
        g_1=
        \left(\frac1{18}X+\frac{11}{81}\right)(18X+28)
        +\frac{16}{81}
    \]
    give the Bézout identity
    \[
        1=
        \left(\frac9{32}X^2-X+\frac{15}{16}\right)g_1
        -
        \left(\frac9{32}X+\frac{11}{16}\right)f.
    \]
    Therefore
    \[
        (g_1+(f))^{-1}
        =
        \left(\frac9{32}X^2-X+\frac{15}{16}\right)+(f).
    \]

    Finally, \(f=0\) in \(A\) gives
    \[
        X^3=2X^2+2X-4.
    \]
    Multiplying by \(X\) and substituting once more,
    \[
        X^4
        =2X^3+2X^2-4X
        =6X^2-8.
    \]
    Hence
    \[
        X^4+(f)=(6X^2-8)+(f).
    \]
\end{snippetsolution}

\end{document}
